\documentclass[11pt]{article}

\usepackage[a4paper,margin=1in]{geometry}
\usepackage{amsmath,amssymb,amsthm,mathtools}
\usepackage{microtype}

\newtheorem{theorem}{Theorem}
\newtheorem{lemma}{Lemma}
\newtheorem{proposition}{Proposition}
\newtheorem{corollary}{Corollary}
\newtheorem{example}{Example}
\newtheorem{remark}{Remark}
\newtheorem{conjecture}{Conjecture}

\DeclareMathOperator{\Tr}{Tr}
\DeclareMathOperator{\Ent}{Ent}
\DeclareMathOperator{\Cov}{Cov}

\newcommand{\R}{\mathbb R}
\newcommand{\cPtwo}{\mathcal P_2}
\newcommand{\cG}{\mathcal G}
\newcommand{\bSplus}{\mathbb S_+}
\newcommand{\W}{W}
\newcommand{\N}{\mathsf N}
\newcommand{\dd}{\,\mathrm d}

\title{A variational criterion for Gaussian-preserving Wasserstein flows}
\author{}
\date{}

\begin{document}

\maketitle

Let
\[
R:\mathcal P_2(\R^d)\to\mathcal P_2(\R^d)
\]
be the map sending a probability measure to the Gaussian measure with the same
mean and covariance:
\[
R\mu:=\N(m_\mu,\Sigma_\mu),
\]
where
\[
m_\mu:=\int_{\R^d}x\,\dd\mu(x),
\qquad
\Sigma_\mu:=\int_{\R^d}(x-m_\mu)(x-m_\mu)^{\mathsf T}\,\dd\mu(x).
\]
The global condition suggested by Hugo Lavenant is
\[
F(R\mu)\le F(\mu)
\qquad
\text{for every }\mu\in\mathcal P_2(\R^d).
\]
This condition is a clean sufficient condition for the JKO scheme, and hence
for the Wasserstein gradient flow when the latter is uniquely obtained as the
JKO limit, to preserve Gaussian measures.

It is important, however, to distinguish this global condition from mere
continuous-time Gaussian preservation. The condition
\[
F(R\mu)\le F(\mu)
\]
is stronger than saying that the Wasserstein gradient flow preserves Gaussian
initial data. A simple smooth cylinder functional gives a counterexample to the
converse.

\section{Gelbrich contraction}

Let
\[
\mathcal G_d
:=
\{\N(m,\Sigma):m\in\R^d,\ \Sigma\in\mathbb S_+^d\}
\]
be the family of possibly degenerate Gaussian laws.

For $A,B\in\mathbb S_+^d$, define
\[
\mathfrak B^2(A,B)
:=
\Tr A+\Tr B
-
2\Tr\!\left((A^{1/2}BA^{1/2})^{1/2}\right).
\]

\begin{theorem}[Gelbrich's inequality]
For every $\mu,\nu\in\mathcal P_2(\R^d)$,
\[
\W_2^2(\mu,\nu)
\ge
\|m_\mu-m_\nu\|^2+\mathfrak B^2(\Sigma_\mu,\Sigma_\nu).
\]
Moreover, for Gaussian laws,
\[
\W_2^2(\N(m,A),\N(n,B))
=
\|m-n\|^2+\mathfrak B^2(A,B).
\]
Consequently,
\[
\W_2(R\mu,R\nu)\le \W_2(\mu,\nu)
\qquad
\text{for every }\mu,\nu\in\mathcal P_2(\R^d).
\]
\end{theorem}

\begin{proof}
The first inequality is Gelbrich's inequality. Applying the Gaussian formula to
$R\mu$ and $R\nu$ gives
\[
\W_2^2(R\mu,R\nu)
=
\|m_\mu-m_\nu\|^2+\mathfrak B^2(\Sigma_\mu,\Sigma_\nu).
\]
Gelbrich's inequality then yields
\[
\W_2^2(R\mu,R\nu)\le \W_2^2(\mu,\nu).
\]
\end{proof}

\section{Lavenant's variational criterion}

Let $F:\mathcal P_2(\R^d)\to(-\infty,+\infty]$ be a functional. For
$\tau>0$ and $\gamma\in\mathcal P_2(\R^d)$, define the JKO objective
\[
\mathcal J_{\tau,\gamma}(\eta)
:=
F(\eta)+\frac{1}{2\tau}\W_2^2(\gamma,\eta).
\]

\begin{theorem}[Lavenant's sufficient condition]
Assume
\[
F(R\mu)\le F(\mu)
\qquad
\text{for every }\mu\in\mathcal P_2(\R^d).
\]
Let $\gamma\in\mathcal G_d$. Then, for every $\eta\in\mathcal P_2(\R^d)$,
\[
\mathcal J_{\tau,\gamma}(R\eta)
\le
\mathcal J_{\tau,\gamma}(\eta).
\]
Consequently, if the JKO problem
\[
\inf_{\eta\in\mathcal P_2(\R^d)}
\left\{
F(\eta)+\frac{1}{2\tau}\W_2^2(\gamma,\eta)
\right\}
\]
admits a minimizer, then it admits a Gaussian minimizer. If the minimizer is
unique, then it is Gaussian.
\end{theorem}

\begin{proof}
Since $\gamma$ is Gaussian,
\[
R\gamma=\gamma.
\]
By Gelbrich contraction,
\[
\W_2(R\eta,\gamma)
=
\W_2(R\eta,R\gamma)
\le
\W_2(\eta,\gamma).
\]
By assumption,
\[
F(R\eta)\le F(\eta).
\]
Therefore
\[
F(R\eta)+\frac{1}{2\tau}\W_2^2(\gamma,R\eta)
\le
F(\eta)+\frac{1}{2\tau}\W_2^2(\gamma,\eta).
\]
Thus projecting any competitor through $R$ can only decrease the JKO objective.

If $\eta$ is a minimizer, then $R\eta$ is also a minimizer and is Gaussian.
If the minimizer is unique, then $\eta=R\eta$, hence $\eta$ is Gaussian.
\end{proof}

\begin{corollary}
Assume
\[
F(R\mu)\le F(\mu)
\qquad
\text{for every }\mu\in\mathcal P_2(\R^d).
\]
If the Wasserstein gradient flow of $F$ is unique and obtained as the limit of
the JKO scheme, then Gaussian initial data remain Gaussian.
\end{corollary}

\section{The exact variational equivalence}

The condition
\[
F(R\mu)\le F(\mu)
\]
is not merely sufficient for the previous proof. It is exactly equivalent to
the fact that Gaussianization improves every JKO objective with Gaussian base
point.

\begin{theorem}
The following two statements are equivalent.

\begin{enumerate}
\item For every $\mu\in\mathcal P_2(\R^d)$,
\[
F(R\mu)\le F(\mu).
\]

\item For every $\tau>0$, every $\gamma\in\mathcal G_d$, and every
$\eta\in\mathcal P_2(\R^d)$,
\[
\mathcal J_{\tau,\gamma}(R\eta)
\le
\mathcal J_{\tau,\gamma}(\eta).
\]
\end{enumerate}
\end{theorem}

\begin{proof}
Assume first that
\[
F(R\mu)\le F(\mu)
\]
for all $\mu$. If $\gamma$ is Gaussian, then $R\gamma=\gamma$, and by Gelbrich
contraction,
\[
\W_2(\gamma,R\eta)
=
\W_2(R\gamma,R\eta)
\le
\W_2(\gamma,\eta).
\]
Therefore
\[
F(R\eta)+\frac{1}{2\tau}\W_2^2(\gamma,R\eta)
\le
F(\eta)+\frac{1}{2\tau}\W_2^2(\gamma,\eta).
\]

Conversely, assume that the JKO comparison holds. Fix
$\eta\in\mathcal P_2(\R^d)$ and choose
\[
\gamma:=R\eta.
\]
Then $\gamma$ is Gaussian, and the comparison gives
\[
F(R\eta)
\le
F(\eta)+\frac{1}{2\tau}\W_2^2(R\eta,\eta).
\]
Letting $\tau\to+\infty$ yields
\[
F(R\eta)\le F(\eta).
\]
\end{proof}

\begin{remark}
Thus the condition
\[
F\circ R\le F
\]
is the right condition for the specific variational argument based on the
retraction $R$. What is not automatic is the converse implication from
continuous-time Gaussian preservation to this global inequality.
\end{remark}

\section{Examples satisfying \(F\circ R\le F\)}

\begin{example}[Functionals of mean and covariance]
Suppose
\[
F(\mu)=\Phi(m_\mu,\Sigma_\mu)
\]
for some map
\[
\Phi:\R^d\times\mathbb S_+^d\to(-\infty,+\infty].
\]
Then
\[
F(R\mu)=F(\mu).
\]
In particular, quadratic interaction energies are included. If
\[
F(\mu)
=
\frac12\int_{\R^d\times\R^d}
(x-y)^{\mathsf T}Q(x-y)\,\dd\mu(x)\dd\mu(y),
\]
with $Q=Q^{\mathsf T}$, then
\[
F(\mu)=\Tr(Q\Sigma_\mu),
\]
so $F$ depends only on the covariance.
\end{example}

\begin{proof}
The identities
\[
m_{R\mu}=m_\mu,
\qquad
\Sigma_{R\mu}=\Sigma_\mu
\]
give the first claim.

For the quadratic interaction, let $X,Y$ be independent random variables with
law $\mu$. Then
\[
\mathbb E[(X-Y)(X-Y)^{\mathsf T}]=2\Sigma_\mu.
\]
Hence
\[
F(\mu)
=
\frac12\Tr\!\left(Q\,\mathbb E[(X-Y)(X-Y)^{\mathsf T}]\right)
=
\Tr(Q\Sigma_\mu).
\]
\end{proof}

\begin{example}[Boltzmann entropy]
Define
\[
\Ent(\mu)
:=
\begin{cases}
\displaystyle \int_{\R^d}\rho\log\rho\,\dd x,
& \mu=\rho\,\dd x,\\[0.8em]
+\infty,
& \text{otherwise}.
\end{cases}
\]
Then
\[
\Ent(R\mu)\le \Ent(\mu).
\]
\end{example}

\begin{proof}
Assume first that $\mu=\rho\,\dd x$ has finite entropy and nondegenerate
covariance. Let $g$ be the density of $R\mu$. The nonnegativity of relative
entropy gives
\[
0\le
\int_{\R^d}\rho\log\frac{\rho}{g}\,\dd x
=
\Ent(\mu)-\int_{\R^d}\rho\log g\,\dd x.
\]
Since $\log g$ is a constant plus a quadratic polynomial, and since $\mu$ and
$R\mu$ have the same mean and covariance,
\[
\int_{\R^d}\rho\log g\,\dd x
=
\int_{\R^d}g\log g\,\dd x
=
\Ent(R\mu).
\]
Therefore
\[
\Ent(R\mu)\le \Ent(\mu).
\]
The degenerate case follows by approximation or by restricting to the affine
span of the covariance.
\end{proof}

\begin{example}[Squared transport cost to a Gaussian]
Let $\gamma\in\mathcal G_d$ and define
\[
F(\mu):=\frac12\W_2^2(\mu,\gamma).
\]
Then
\[
F(R\mu)\le F(\mu).
\]
\end{example}

\begin{proof}
Since $\gamma$ is Gaussian,
\[
R\gamma=\gamma.
\]
By Gelbrich contraction,
\[
\W_2(R\mu,\gamma)
=
\W_2(R\mu,R\gamma)
\le
\W_2(\mu,\gamma).
\]
Squaring gives the claim.
\end{proof}

\section{Why the converse is false without extra assumptions}

One may ask whether Gaussian preservation of the continuous-time Wasserstein
gradient flow implies
\[
F(R\mu)\le F(\mu)
\qquad
\text{for all }\mu.
\]
In this generality, the answer is no. The reason is that flow preservation is a
local condition along the Gaussian manifold, whereas
\[
F(R\mu)\le F(\mu)
\]
is a global comparison between a measure and its Gaussianization.

\begin{proposition}[A counterexample to necessity]
There exists a smooth bounded functional $F:\mathcal P_2(\R)\to(0,1]$ such
that:

\begin{enumerate}
\item every Gaussian measure is a stationary point of the formal Wasserstein
gradient flow of $F$;

\item hence, under uniqueness of the flow, Gaussian initial data remain
Gaussian;

\item nevertheless, there exists $\mu\in\mathcal P_2(\R)$ such that
\[
F(R\mu)>F(\mu).
\]
\end{enumerate}
Thus Gaussian preservation of the flow does not imply
\[
F\circ R\le F.
\]
\end{proposition}

\begin{proof}
For $\mu\in\mathcal P_2(\R)$, write
\[
m_\mu:=\int x\,\dd\mu(x),
\qquad
\sigma_\mu^2:=\int (x-m_\mu)^2\,\dd\mu(x).
\]
Define
\[
D(\mu)
:=
\int_{\R}\cos x\,\dd\mu(x)
-
e^{-\sigma_\mu^2/2}\cos(m_\mu),
\]
and set
\[
F(\mu):=\exp(-D(\mu)^2).
\]

If $\mu=\N(m,\sigma^2)$, then
\[
\int_{\R}\cos x\,\dd\mu(x)
=
e^{-\sigma^2/2}\cos m.
\]
Therefore
\[
D(\mu)=0
\qquad
\text{for every Gaussian }\mu,
\]
and hence
\[
F(\mu)=1
\qquad
\text{for every Gaussian }\mu.
\]

We next compute the Wasserstein differential. Let $(\mu_t)_{t}$ be a smooth
curve solving
\[
\partial_t\mu_t+\partial_x(\mu_t v_t)=0.
\]
Then
\[
\frac{\dd}{\dd t}m_{\mu_t}
=
\int v_t\,\dd\mu_t,
\]
and
\[
\frac{\dd}{\dd t}\sigma_{\mu_t}^2
=
2\int (x-m_{\mu_t})v_t(x)\,\dd\mu_t(x).
\]
Also,
\[
\frac{\dd}{\dd t}\int\cos x\,\dd\mu_t(x)
=
\int -\sin x\,v_t(x)\,\dd\mu_t(x).
\]
Thus
\[
\frac{\dd}{\dd t}D(\mu_t)
=
\int G_{\mu_t}(x)v_t(x)\,\dd\mu_t(x),
\]
where
\[
G_\mu(x)
=
-\sin x
+
e^{-\sigma_\mu^2/2}
\left[
\sin(m_\mu)+\cos(m_\mu)(x-m_\mu)
\right].
\]
Consequently,
\[
\nabla_W D(\mu)(x)=G_\mu(x),
\]
and
\[
\nabla_WF(\mu)(x)
=
-2D(\mu)e^{-D(\mu)^2}G_\mu(x).
\]
If $\gamma$ is Gaussian, then $D(\gamma)=0$, so
\[
\nabla_WF(\gamma)=0.
\]
Therefore every Gaussian is a stationary solution of the formal Wasserstein
gradient-flow equation
\[
\partial_t\mu_t
=
\partial_x\left(\mu_t\nabla_WF(\mu_t)\right).
\]
Under uniqueness, the flow starting from a Gaussian remains at that Gaussian.

It remains to check that the inequality
\[
F(R\mu)\le F(\mu)
\]
fails. Take
\[
\mu:=\frac12\delta_{-1}+\frac12\delta_1.
\]
Then
\[
m_\mu=0,
\qquad
\sigma_\mu^2=1,
\qquad
R\mu=\N(0,1).
\]
Hence
\[
D(R\mu)=0,
\]
while
\[
D(\mu)
=
\cos(1)-e^{-1/2}
\neq 0.
\]
Therefore
\[
F(R\mu)=1
\]
and
\[
F(\mu)
=
\exp\left(-(\cos(1)-e^{-1/2})^2\right)
<1.
\]
Thus
\[
F(R\mu)>F(\mu),
\]
which is the opposite of the desired inequality.
\end{proof}

\begin{remark}
This counterexample is intentionally simple. It shows that a functional may have
zero Wasserstein gradient on the whole Gaussian manifold, so that Gaussian
measures are stationary, while assigning larger values to Gaussian measures
than to some non-Gaussian measures with the same mean and covariance.
\end{remark}

\section{Conclusion}

The condition
\[
F(R\mu)\le F(\mu)
\]
has a precise variational meaning: it is equivalent to saying that the
Gaussianization map $R$ improves every JKO objective with Gaussian base point.
Together with Gelbrich's contraction, this gives a direct proof that the JKO
scheme can be selected inside the Gaussian manifold.

However, this condition is not necessary for the weaker property that the
continuous-time Wasserstein gradient flow preserves Gaussian initial data.
Preservation of the flow is local along $\mathcal G_d$, whereas
\[
F(R\mu)\le F(\mu)
\]
is a global comparison between $\mu$ and its moment-matched Gaussian.

\begin{thebibliography}{9}

\bibitem{Gelbrich1990}
M.~Gelbrich,
\newblock On a formula for the $L^2$-Wasserstein metric between measures on
Euclidean and Hilbert spaces,
\newblock \emph{Mathematische Nachrichten} \textbf{147} (1990), 185--203.

\bibitem{LavenantPC}
H.~Lavenant,
\newblock Gaussian-preserving Wasserstein flows,
\newblock personal communication, October 2025.

\end{thebibliography}

\end{document}